% Gerado de tcc.md (exportado do Claude Docs) por ferramentas/gera_latex.php em 12/11/2026 21:45.
% Para mudar o texto, edite este arquivo; se gerar de novo a partir do documento, as edições feitas aqui se perdem.

\capitulosemnumero{REFERÊNCIAS}

\begin{referenciasIFTO}
ALMEIDA, Fernanda Correia de Lima. \textbf{Plano educacional individualizado}: aplicativo de uso colaborativo na elaboração do PEI para um aluno autista. 2023. 115 f. Dissertação (Mestrado em Educação) – Programa de Pós-Graduação em Educação, Universidade Federal do Tocantins, Palmas, 2023. Disponível em: \url{https://repositorio.uft.edu.br/handle/11612/6107}. Acesso em: 12 nov. 2026.

BENTES, Juliana Ferreira; ARAÚJO, Fabiola Pantoja Oliveira; PIRES, Yomara Pinheiro. ANE - jogos educativos para auxiliar crianças com deficiência intelectual no aprendizado do alfabeto, numerais e emoções. \textit{In}: WORKSHOP DE INFORMÁTICA NA ESCOLA, 26., 2020, On-line. \textbf{Anais} [...]. Porto Alegre: Sociedade Brasileira de Computação, 2020. p. 399-408. DOI: \url{https://doi.org/10.5753/cbie.wie.2020.399}.

BRASIL. [Constituição (1988)]. \textbf{Constituição da República Federativa do Brasil de 1988}. Brasília, DF: Presidência da República, [2026]. Disponível em: \url{https://www.planalto.gov.br/ccivil_03/constituicao/constituicao.htm}. Acesso em: 9 nov. 2026.

BRASIL. Lei nº 9.394, de 20 de dezembro de 1996. Estabelece as diretrizes e bases da educação nacional. \textbf{Diário Oficial da União}: seção 1, Brasília, DF, 23 dez. 1996. Disponível em: \url{https://www.planalto.gov.br/ccivil_03/leis/l9394.htm}. Acesso em: 9 nov. 2026.

BRASIL. Ministério da Educação. Secretaria de Educação Especial. \textbf{Política Nacional de Educação Especial na Perspectiva da Educação Inclusiva}. Brasília, DF: MEC/SEESP, 2008.

BRASIL. Ministério da Educação. Conselho Nacional de Educação. Câmara de Educação Básica. Resolução nº 4, de 2 de outubro de 2009. Institui Diretrizes Operacionais para o Atendimento Educacional Especializado na Educação Básica, modalidade Educação Especial. \textbf{Diário Oficial da União}: seção 1, Brasília, DF, 5 out. 2009.

BRASIL. Decreto nº 7.611, de 17 de novembro de 2011. Dispõe sobre a educação especial, o atendimento educacional especializado e dá outras providências. \textbf{Diário Oficial da União}: seção 1, Brasília, DF, 18 nov. 2011.

BRASIL. Lei nº 13.146, de 6 de julho de 2015. Institui a Lei Brasileira de Inclusão da Pessoa com Deficiência (Estatuto da Pessoa com Deficiência). \textbf{Diário Oficial da União}: seção 1, Brasília, DF, 7 jul. 2015. Disponível em: \url{https://www.planalto.gov.br/ccivil_03/_ato2015-2018/2015/lei/l13146.htm}. Acesso em: 9 nov. 2026.

BRASIL. Lei nº 13.709, de 14 de agosto de 2018. Lei Geral de Proteção de Dados Pessoais (LGPD). \textbf{Diário Oficial da União}: seção 1, Brasília, DF, 15 ago. 2018a.

BRASIL. Ministério da Educação. \textbf{Base Nacional Comum Curricular}. Brasília, DF: MEC, 2018b. Disponível em: \url{http://basenacionalcomum.mec.gov.br/images/BNCC_EI_EF_110518_versaofinal_site.pdf}. Acesso em: 9 nov. 2026.

BRASIL. Decreto nº 12.686, de 20 de outubro de 2025. Institui a Política Nacional de Educação Especial Inclusiva e a Rede Nacional de Educação Especial Inclusiva. \textbf{Diário Oficial da União}: seção 1, Brasília, DF, 21 out. 2025a. Disponível em: \url{https://www.planalto.gov.br/ccivil_03/_ato2023-2026/2025/decreto/d12686.htm}. Acesso em: 9 nov. 2026.

BRASIL. Decreto nº 12.773, de 8 de dezembro de 2025. Altera o Decreto nº 12.686, de 20 de outubro de 2025, que institui a Política Nacional de Educação Especial Inclusiva e a Rede Nacional de Educação Especial Inclusiva. \textbf{Diário Oficial da União}: seção 1, Brasília, DF, 9 dez. 2025b. Disponível em: \url{https://www.planalto.gov.br/ccivil_03/_ato2023-2026/2025/decreto/d12773.htm}. Acesso em: 9 nov. 2026.

BRASIL. Ministério da Educação. Portaria nº 421, de 15 de maio de 2026. Dispõe sobre a Política Nacional de Educação Especial Inclusiva - PNEEI e sobre a Rede Nacional de Educação Especial Inclusiva - Reneei, de que tratam o Decreto nº 12.686, de 20 de outubro de 2025. \textbf{Diário Oficial da União}: Brasília, DF, 18 maio 2026. Texto consolidado, com as alterações da Portaria MEC nº 550, de 18 de junho de 2026.

DENO, Stanley L. Curriculum-based measurement: the emerging alternative. \textbf{Exceptional Children}, [\textit{s. l.}], v. 52, n. 3, p. 219-232, 1985. DOI: \url{https://doi.org/10.1177/001440298505200303}.

DINIZ, Daniele C. S. \textit{et al.} PAR Digital: uma ferramenta web em prol da educação inclusiva. \textit{In}: SIMPÓSIO BRASILEIRO DE SISTEMAS MULTIMÍDIA E WEB, 30., 2024, Juiz de Fora. \textbf{Anais Estendidos} [...]. Porto Alegre: Sociedade Brasileira de Computação, 2024. p. 107-110. DOI: \url{https://doi.org/10.5753/webmedia_estendido.2024.244113}.

DRESCH, Aline; LACERDA, Daniel Pacheco; ANTUNES JÚNIOR, José Antonio Valle. \textbf{Design science research}: método de pesquisa para avanço da ciência e tecnologia. Porto Alegre: Bookman, 2015.

DUTRA, Taynara Cerigueli; MASCHIO, Eleandro; GASPARINI, Isabela. Pensar e Lavar: processo de desenvolvimento e avaliação de um jogo digital educacional para promover o pensamento computacional para crianças neurotípicas e com deficiência intelectual. \textbf{Revista Brasileira de Informática na Educação}, Porto Alegre, v. 31, p. 659-690, 2023. DOI: \url{https://doi.org/10.5753/rbie.2023.3330}.

FARIAS, Ezequiel B.; SILVA, Leandro W. C.; CUNHA, Mônica X. C. ABC Autismo: um aplicativo móvel para auxiliar na alfabetização de crianças com autismo baseado no Programa TEACCH. \textit{In}: SIMPÓSIO BRASILEIRO DE SISTEMAS DE INFORMAÇÃO, 10., 2014, Londrina. \textbf{Anais} [...]. Porto Alegre: Sociedade Brasileira de Computação, 2014. p. 458-469. DOI: \url{https://doi.org/10.5753/sbsi.2014.6136}.

FERNANDES, Marcela Peron. \textbf{Estratégias de gamificação no processo de alfabetização de alunos com Transtorno do Espectro Autista (TEA)}. 2026. 133 f. Dissertação (Mestrado Profissional em Educação Inclusiva) – Faculdade de Ciências e Tecnologia, Universidade Estadual Paulista, Presidente Prudente, 2026. Disponível em: \url{https://hdl.handle.net/11449/320939}. Acesso em: 9 nov. 2026.

FOWLER, Martin. \textbf{Padrões de arquitetura de aplicações corporativas}. Porto Alegre: Bookman, 2006.

FRANCISCO, Priscila Rodrigues Cruz. \textbf{Elaboração do Plano Educacional Individualizado (PEI)}: proposta de um recurso educacional digital. 2023. 146 f. Dissertação (Mestrado Profissional em Ensino na Educação Básica) – Instituto de Aplicação Fernando Rodrigues da Silveira, Universidade do Estado do Rio de Janeiro, Rio de Janeiro, 2023. Disponível em: \url{https://www.bdtd.uerj.br:8443/bitstream/1/20629/2/Disserta%c3%a7%c3%a3o%20-%20Priscila%20Rodrigues%20Cruz%20Francisco%20-%20Completa%20-%202023.pdf}. Acesso em: 9 nov. 2026.

GALDINO, Jaqueline Leite Silva. \textbf{Alfabetização matemática de alunos com TEA}: uma proposta de ensino mediada pela tecnologia. 2024. 77 f. Dissertação (Mestrado em Matemática) – Departamento de Matemática, Pontifícia Universidade Católica do Rio de Janeiro, Rio de Janeiro, 2024. Disponível em: \url{https://www.maxwell.vrac.puc-rio.br/66367/66367.PDF}. Acesso em: 9 nov. 2026.

GIL, Antonio Carlos. \textbf{Como elaborar projetos de pesquisa}. 6. ed. São Paulo: Atlas, 2017.

INSTITUTO NACIONAL DE ESTUDOS E PESQUISAS EDUCACIONAIS ANÍSIO TEIXEIRA (INEP). \textbf{Censo da Educação Básica 2020}: notas estatísticas. Brasília, DF: Inep, 2021. Disponível em: \url{https://download.inep.gov.br/publicacoes/institucionais/estatisticas_e_indicadores/notas_estatisticas_censo_escolar_2020.pdf}. Acesso em: 12 nov. 2026.

INSTITUTO NACIONAL DE ESTUDOS E PESQUISAS EDUCACIONAIS ANÍSIO TEIXEIRA (INEP). \textbf{Censo Escolar da Educação Básica 2024}: notas estatísticas. Brasília, DF: Inep, 2025. Disponível em: \url{https://download.inep.gov.br/publicacoes/institucionais/estatisticas_e_indicadores/notas_estatisticas_censo_da_educacao_basica_2024.pdf}. Acesso em: 12 nov. 2026.

INSTITUTO NACIONAL DE ESTUDOS E PESQUISAS EDUCACIONAIS ANÍSIO TEIXEIRA (INEP). \textbf{Censo Escolar da Educação Básica 2025}: notas estatísticas. Brasília, DF: Inep, 2026. Disponível em: \url{https://download.inep.gov.br/publicacoes/institucionais/estatisticas_e_indicadores/notas_estatisticas_censo_escolar_da_educacao_basica_2025.pdf}. Acesso em: 12 nov. 2026.

MESQUITA, Jeferson Antônio Gonçalves de; BASTOS, Lorena Regina Cruz. Garimpeirinhos: ecossistema digital integrado para alfabetização, inclusão e apoio ao professor. \textit{In}: CONGRESSO BRASILEIRO DE INFORMÁTICA NA EDUCAÇÃO, 14., 2025, Curitiba. \textbf{Anais Estendidos} [...]. Porto Alegre: Sociedade Brasileira de Computação, 2025. p. 200-203. DOI: \url{https://doi.org/10.5753/cbie_estendido.2025.13608}.

MOZILLA. Web Speech API. \textit{In}: MDN WEB DOCS. [\textit{S. l.}]: Mozilla, 2026. Disponível em: \url{https://developer.mozilla.org/en-US/docs/Web/API/Web_Speech_API}. Acesso em: 9 nov. 2026.

NEGRIM, Márcia Regina Corrêa. \textbf{Gestão escolar, docência e tecnologia digital}: trabalho colaborativo para o ensino de alunos com deficiência intelectual. 2019. Dissertação (Mestrado em Docência para a Educação Básica) – Faculdade de Ciências, Universidade Estadual Paulista, Bauru, 2019. Disponível em: \url{https://repositorio.unesp.br/handle/11449/183670}. Acesso em: 12 nov. 2026.

NIELSEN, Jakob. Enhancing the explanatory power of usability heuristics. \textit{In}: SIGCHI CONFERENCE ON HUMAN FACTORS IN COMPUTING SYSTEMS, 1994, Boston. \textbf{Proceedings} [...]. New York: Association for Computing Machinery, 1994a. p. 152-158. DOI: \url{https://doi.org/10.1145/191666.191729}.

NIELSEN, Jakob. \textbf{Severity ratings for usability problems}. [\textit{S. l.}]: Nielsen Norman Group, 1 nov. 1994b. Disponível em: \url{https://www.nngroup.com/articles/how-to-rate-the-severity-of-usability-problems/}. Acesso em: 9 nov. 2026.

NIELSEN, Jakob; MOLICH, Rolf. Heuristic evaluation of user interfaces. \textit{In}: SIGCHI CONFERENCE ON HUMAN FACTORS IN COMPUTING SYSTEMS, 1990, Seattle. \textbf{Proceedings} [...]. New York: Association for Computing Machinery, 1990. p. 249-256. DOI: \url{https://doi.org/10.1145/97243.97281}.

OBJECT MANAGEMENT GROUP. \textbf{OMG Unified Modeling Language (OMG UML)}: version 2.5.1. [\textit{S. l.}]: OMG, dez. 2017. Documento formal/2017-12-05. Disponível em: \url{https://www.omg.org/spec/UML/2.5.1/PDF}. Acesso em: 11 nov. 2026.

OWASP FOUNDATION. \textbf{OWASP Top 10:2021}. [\textit{S. l.}]: OWASP Foundation, 2021. Disponível em: \url{https://owasp.org/Top10/}. Acesso em: 9 nov. 2026.

PAIVA, Willian Micael de; MUNIZ, Raphael de Carvalho; MOREIRA, Irlan Arley Targino; SOUSA, Reudismam Rolim de. Desenvolvimento de um sistema digital para gestão de Planos Educacionais Individualizados (PEIs): promovendo a educação inclusiva por meio de tecnologias. \textbf{Temática}, João Pessoa, ano 22, n. 2, p. 58-82, fev. 2026. Disponível em: \url{https://periodicos.ufpb.br/index.php/tematica/article/view/77906}. Acesso em: 11 nov. 2026.

PRENSKY, Marc. \textbf{Aprendizagem baseada em jogos digitais}. São Paulo: Editora Senac São Paulo, 2012.

PRODANOV, Cleber Cristiano; FREITAS, Ernani Cesar de. \textbf{Metodologia do trabalho científico}: métodos e técnicas da pesquisa e do trabalho acadêmico. 2. ed. Novo Hamburgo: Feevale, 2013.

SILVA, Arthur H. S.; SANTOS, Pedro H.; CUNHA, Mônica X. C. LetraMundo: um jogo sério para alfabetização de crianças com Transtorno do Espectro Autista (TEA) em ambiente escolar. \textit{In}: SIMPÓSIO BRASILEIRO DE INFORMÁTICA NA EDUCAÇÃO, 35., 2024, Rio de Janeiro. \textbf{Anais} [...]. Porto Alegre: Sociedade Brasileira de Computação, 2024. p. 2747-2756. DOI: \url{https://doi.org/10.5753/sbie.2024.244666}.

SOMMERVILLE, Ian. \textbf{Engenharia de software}. 9. ed. São Paulo: Pearson Prentice Hall, 2011.

TANNÚS-VALADÃO, Gabriela; MENDES, Enicéia Gonçalves. Inclusão escolar e o planejamento educacional individualizado: estudo comparativo sobre práticas de planejamento em diferentes países. \textbf{Revista Brasileira de Educação}, Rio de Janeiro, v. 23, e230076, p. 1-18, 2018. DOI: \url{https://doi.org/10.1590/S1413-24782018230076}.

VIGOTSKI, Lev Semenovich. \textbf{A formação social da mente}: o desenvolvimento dos processos psicológicos superiores. 7. ed. São Paulo: Martins Fontes, 2007.

WATT, Adrienne; ENG, Nelson. \textbf{Database design}. 2. ed. Victoria, B.C.: BCcampus, 2014. Disponível em: \url{https://dc.arcabc.ca/_flysystem/repo-bin/2021-10/dc_51320_1.pdf}. Acesso em: 11 nov. 2026.

WORLD WIDE WEB CONSORTIUM (W3C). \textbf{Web Content Accessibility Guidelines (WCAG) 2.1}. [\textit{S. l.}]: W3C, 2018. Disponível em: \url{https://www.w3.org/TR/WCAG21/}. Acesso em: 9 nov. 2026.

\end{referenciasIFTO}
